% page 287 of the facsimile (image p-286 of the scan)
% block 167.6 x 254.9 mm ; left margin 34.4 mm
\pgc{12.700mm}{0.000mm}{
\l{}
\l{}
\l{}
\l{}
\l{\cel{53}279}
\l{}
\l{\sou{klematido}. (\sou{bot.}) Planto klimera, familio \textquotedbl{}ranunkulacei\textquotedbl{}, di}
\l{\cel{3}qua la flori esas blanka ed odoroza. - L.\cel{1}\sou{clematis}. EFIS.}
\l{}
\l{\sou{klementa}. Dicesas pri persono qua, autoritatoza pri punisar kul-}
\l{\cel{3}pozo, agas dolce ad ica, pardonante lu, od alejante lua puni-}
\l{\cel{3}seso. - EFIS.}
\l{}
\l{\sou{klepsidro}. Aquo-horlojo qua indikas la kloki per la fluo ne-inter-}
\l{\cel{3}ruptita da ula quanto de aquo en tempo-quanto predeterminita. -}
\l{\cel{3}- DEFIS.}
\l{}
\l{\sou{kleriko}. Membro di la oficieraro di eklezio. - DEFIRS.}
\l{}
\l{\sou{klerko}. La persono qua laboras o yuro-studias en la kontoro di}
\l{\cel{3}notario, di ushero, di \textquotedbl{}avoué\textquotedbl{}. - EF.}
\l{}
\l{\sou{kliento.}\cel{2}I. (\sou{Roma, antiqua}). Plebeyo qua protektesis da patrico.}
\l{\cel{3}- II. La persono qua konfidas kustumale sua interesti ad afer-}
\l{\cel{3}isto (advokato, notario, e c.), sua saneso a mediko ; la per-}
\l{\cel{3}sono di qua advokato pledas la afero. - III. (komerco). La per-}
\l{\cel{3}sono qua kompras kustumale de komercisto, qua frequentas butiko}
\l{\cel{3}(restorerio, e c.). - DEFIRS.}
\l{}
\l{\sou{klifo}. Eskarpajo de tero o de roko qua garnisas la bordo di}
\l{\cel{3}maro. - E.}
\l{}
\l{\sou{kliko}. (\sou{tekn}.) Peco movebla quan sublevas omna dento di ingrano}
\l{\cel{3}roto kande ica turnas ica-sinse, e qua butas kontre la ingran}
\l{\cel{3}ajo se onu volas igar lu turnar inversa-sinse. - DEF.}
\l{}
\l{\sou{kliktar}. (\sou{netrans}.)(\sou{aludante ula korpi sonora qui intershokas).}}
\l{\cel{3}Bruisar ne-matide. - DEF.}
\l{}
\l{klimar. (\sou{netrans}.)Elevar su per akrochar su a to quo povas helpar}
\l{\cel{3}la acenso. - DE.}
\l{}
\l{klimato. Ensemblo ek la cirkonstanci atmosferala, a qua regioni}
\l{\cel{3}submisesas. - DEFIRS.}
\l{}
\l{klimatologio. Studio di la klimati. - DEFS.}
\l{}
\l{klinar.\cel{2}(\sou{netrans.)}\cel{1}(\sou{aludante du peci qui quaze interkavalkas).}}
\l{\cel{3}Interkrucumar. - F.}
\l{}
\l{kliniko. Docado medicinala, medikala, da maestro a sua dicipuli,}
\l{\cel{3}an la lito di la maladi. - DEFIRS.}
\l{}
\l{\cel{1}\sou{klinko}. Klozilo di pordo, ek simpla fero-lamo quan onu abasas}
\l{\cel{4}sur fero-peco muntita an la pordo-kadro, kande onu klozas}
\l{\cel{4}la pordo e sublevas kande onu apertas lu. - DF.}
\l{}
\l{\cel{1}klinoida. (\sou{anat.}) Singla ek la quar apofizi, qui esas an la parto}
\l{\cel{4}supra di la osto sfenoida.}
}
